\diarytitle{0069}{波動方程式（両端固定）における三角波状初期変位のフーリエ級数解}

長さ $L$ の弦の両端を固定した場合の初期値・境界値問題は、変数分離法（フーリエ級数展開）により解く。変数分離 $u=X(x)T(t)$ と境界条件から固有値 $\lambda_n=(n\pi/L)^2$、固有関数 $X_n(x)=\sin(n\pi x/L)$ が定まり、一般解は
\[
u(x,t) = \sum_{n=1}^{\infty}\left(A_{n}\cos\frac{n\pi c t}{L} + B_{n}\sin\frac{n\pi c t}{L}\right)\sin\frac{n\pi x}{L}
\]
となる（導出は本テーマ他問と同様）。初期条件に代入すると
\[
u(x,0) = \sum_{n=1}^{\infty} A_{n}\sin\frac{n\pi x}{L} = f(x), \qquad
\frac{\partial u}{\partial t}(x,0) = \sum_{n=1}^{\infty} \frac{n\pi c}{L}B_{n}\sin\frac{n\pi x}{L} = 0
\]
であり、直交性 $\displaystyle\int_{0}^{L}\sin\frac{n\pi x}{L}\sin\frac{m\pi x}{L}\,dx = \frac{L}{2}\delta_{nm}$ から係数が定まる。

\medskip
\noindent
\textbf{初期速度の項}\quad $u_{t}(x,0)=0$ より、すべての $n$ について $B_{n}=0$。

\medskip
\noindent
\textbf{初期変位の項（$A_n$ の計算）}\quad
\[
A_{n} = \frac{2}{L}\int_{0}^{L} f(x)\sin\frac{n\pi x}{L}\,dx
= \frac{2}{L}\left[\int_{0}^{L/2}\frac{2h}{L}x\sin\frac{n\pi x}{L}\,dx + \int_{L/2}^{L}\frac{2h}{L}(L-x)\sin\frac{n\pi x}{L}\,dx\right]
\]
第2項は $x' = L-x$ と置換すると
\[
\int_{0}^{L/2}\frac{2h}{L}x'\sin\frac{n\pi(L-x')}{L}\,dx'
= \int_{0}^{L/2}\frac{2h}{L}x'\left(\sin n\pi\cos\frac{n\pi x'}{L} - \cos n\pi\sin\frac{n\pi x'}{L}\right)dx'
\]
となる。$\sin n\pi = 0$、$\cos n\pi = (-1)^{n}$ であるから
\[
= -(-1)^{n}\int_{0}^{L/2}\frac{2h}{L}x'\sin\frac{n\pi x'}{L}\,dx'
\]
よって
\[
A_{n} = \frac{2}{L}\cdot\frac{2h}{L}\left(1-(-1)^{n}\right)\int_{0}^{L/2} x\sin\frac{n\pi x}{L}\,dx
\]
偶数 $n$ では $1-(-1)^n=0$ より $A_{n}=0$。奇数 $n$ では $1-(-1)^{n}=2$ であり、部分積分
\[
\int_{0}^{L/2} x\sin\frac{n\pi x}{L}\,dx
= \left[-\frac{L}{n\pi}x\cos\frac{n\pi x}{L}\right]_{0}^{L/2} + \frac{L}{n\pi}\int_{0}^{L/2}\cos\frac{n\pi x}{L}\,dx
= -\frac{L^{2}}{2n\pi}\cos\frac{n\pi}{2} + \left(\frac{L}{n\pi}\right)^{2}\sin\frac{n\pi}{2}
\]
奇数 $n$ では $\cos(n\pi/2)=0$ なので
\[
\int_{0}^{L/2} x\sin\frac{n\pi x}{L}\,dx = \left(\frac{L}{n\pi}\right)^{2}\sin\frac{n\pi}{2}
\]
これらをまとめると
\[
A_{n} = \frac{2}{L}\cdot\frac{2h}{L}\cdot 2 \cdot \left(\frac{L}{n\pi}\right)^{2}\sin\frac{n\pi}{2} = \frac{8h}{n^{2}\pi^{2}}\sin\frac{n\pi}{2}
\qquad (n:\text{奇数}),\qquad A_{n}=0\ (n:\text{偶数})
\]
$\sin(n\pi/2)$ は $n=1,3,5,7,\dots$ に対して $1,-1,1,-1,\dots$ となるので、$n=2m-1$（$m=1,2,3,\dots$）と書けば $\sin\frac{n\pi}{2} = (-1)^{m-1}$。よって
\[
\boxed{\;
u(x,t) = \frac{8h}{\pi^{2}}\sum_{m=1}^{\infty}\frac{(-1)^{m-1}}{(2m-1)^{2}}\cos\frac{(2m-1)\pi c t}{L}\,\sin\frac{(2m-1)\pi x}{L}
\;}
\]

（検算: $t=0$ とすると $u(x,0) = \dfrac{8h}{\pi^{2}}\sum_{m=1}^{\infty}\dfrac{(-1)^{m-1}}{(2m-1)^{2}}\sin\dfrac{(2m-1)\pi x}{L}$ となり、これは三角波 $f(x)$ の標準的なフーリエ正弦級数展開に一致する。特に $x=L/2$ で $\sin\frac{(2m-1)\pi}{2}=(-1)^{m-1}$ なので $u(L/2,0) = \dfrac{8h}{\pi^2}\sum_{m=1}^\infty \dfrac{1}{(2m-1)^2} = \dfrac{8h}{\pi^2}\cdot\dfrac{\pi^2}{8} = h$ となり $f(L/2)=h$ と一致する。また各項は波動方程式・境界条件・$u_t(x,0)=0$ を満たす。）
